\documentclass[11pt]{article}
\usepackage[margin=0.9in]{geometry}
\usepackage{fontspec}\setmainfont{TeX Gyre Termes}
\usepackage{tabularx,booktabs}\usepackage{enumitem}\setlist{nosep,leftmargin=1.2em}
\usepackage{titlesec}\titleformat{\section}{\normalfont\bfseries}{\thesection.}{0.5em}{}\titlespacing*{\section}{0pt}{9pt}{3pt}
\pagestyle{plain}\setlength{\parindent}{0pt}\setlength{\parskip}{4pt}
\begin{document}
{\large\bfseries Progress Report — ACCESS project AGR260032}\\[2pt]
{\small Leakage-controlled, evidence-aware gene annotation for plant genomes. PI: Won Yim, University of Nevada, Reno. Resource: NCSA DeltaAI (GH200). Report date: October 3, 2026.}

\section{Project goal}
We are training a deep-learning gene-structure annotation model for plant genomes. A model of about 400 million parameters reads up to 129,024 nt of genomic sequence and outputs gene structures. A sequence-only baseline is compared with a new model that also reads evidence tracks from ESTs, PacBio and ONT long reads, and cross-species protein alignments. Every evidence record is filtered so that no held-out locus used for validation or testing can leak into training inputs. The pre-registered question is whether the evidence improves gene prediction beyond the sequence-only model, tested with control experiments that isolate what part of the evidence drives any gain.

\section{Use of the current allocation}
Of 1,499 GPU-hours awarded, about 650 have been used and 849 remain.

\begin{tabularx}{\linewidth}{Xrl}\toprule
Activity & GPU-h & Outcome\\\midrule
Sequence-only baseline, seed 123, 4 GPUs, NGC PyTorch 26.08 with torch.compile & 630 & Epoch 13; validation loss still improving; early stopping running\\
Runtime qualification over three sequence-length tiers, 1--4 GPUs & 13 & Production runtime selected; time and memory measured per tier\\
Evidence-aware model integration tests & 5 & Memory gate passed at 63 GB per GPU (limit 100 GB); byte-identical checkpoint resume; distributed fault handling verified\\\bottomrule
\end{tabularx}

\section{Preparatory work on campus resources (no ACCESS charge)}
\begin{itemize}
\item Built 28 leakage-filtered evidence databases for 9 plant species from ESTs, PacBio, ONT and protein alignments. Each species passed exhaustiveness of at least 0.999, zero leakage and all integrity checks.
\item Merged them into one 786 GB evidence database and re-validated the training corpus against it; 37.85 million molecules were consistent with the leakage repair.
\item Built a 654-million-row training cache in 455 shards. Every input is bound by cryptographic hash and re-verified at each job start.
\item All inputs are staged on DeltaAI, and the first evidence-aware training run is ready to submit.
\end{itemize}

\section{Remaining work and GPU-hour requirement}
Measured on DeltaAI with 4 GPUs: about 45 GPU-hours per epoch for the baseline model and about 49.5 for the evidence-aware model. Early stopping (patience 3, maximum 22 epochs) is expected near epoch 18. Our protocol requires the confirmatory baseline seeds 456 and 789 to finish before submission.

\begin{tabularx}{\linewidth}{Xrr}\toprule
Work & Expected & Maximum\\\midrule
Finish baseline seed 123 & 270 & 411\\
Baseline confirmatory seed 456 & 812 & 992\\
Baseline confirmatory seed 789 & 812 & 992\\
Fixed-budget pilot: evidence-aware model and sequence-only counterpart on the same data (4,099 updates each) & 196 & 196\\
Full training of both models to early stopping & 1,584 & 1,980\\
Three fixed-budget controls: shuffled evidence, presence-only evidence, reading-frame channels removed & 294 & 294\\
Evaluation on held-out panels in 9 species and an independent maize test & 50 & 50\\\midrule
Total needed & 4,018 & 4,915\\
Remaining balance & 849 & 849\\
Shortfall & 3,169 & 4,066\\
Contingency (15\%) for queue preemption and one re-run & 603 & 737\\\midrule
\textbf{Supplement requested} & \textbf{3,800} & \\\bottomrule
\end{tabularx}

\section{Expected outcomes}
A pre-registered, leakage-controlled test of whether transcript and protein evidence improves deep-learning gene annotation in plants, with three random seeds for the baseline and control experiments that attribute any gain to evidence content, evidence detail and reading-frame information. Code, frozen protocols and evaluation panels will be released with the manuscript.
\end{document}
